
\section{Conclusión asociada al objetivo general}

Este trabajo desarrolló y validó una metodología basada en ciencia de datos para pronosticar y explicar los indicadores SAIDI y SAIFI en la red de distribución de energía eléctrica de una empresa del Norte de Santander, a partir de la integración de seis fuentes de datos heterogéneas (BRAE, OMS SPARD/SP7, GIS, NASA POWER, más sismicidad, desastres y vegetación) en un panel único cuadrante--semana. La metodología combinó modelos de pronóstico semanal (Pooled OLS, efectos fijos y Random Forest) con modelos espaciales formales (SAR, SEM), obteniendo en ambos casos desempeños satisfactorios para sus respectivos propósitos: predicción bajo incertidumbre temporal real y explicación estructural del patrón territorial de calidad del servicio. El objetivo general se cumple en la medida en que el pipeline completo, desde la extracción de datos crudos hasta los modelos validados, quedó documentado de forma reproducible y aplicable por la empresa a su labor de gestión diaria de la red.

\section{Conclusión asociada al objetivo específico 1}

Se consolidó un repositorio analítico unificado y trazable a partir de seis fuentes complementarias, documentado mediante una tabla de linaje de datos (Tabla~\ref{tab:data_lineage}) que especifica, para cada fuente, las variables extraídas, el proceso de depuración aplicado y el resultado obtenido. El panel final no presenta valores faltantes ($0$~\% sobre $N_{\mathrm{agr}}=90{.}128$ observaciones cuadrante--semana), y el proceso de integración permitió identificar y documentar explícitamente una limitación de cobertura —un sesgo sistemático de emparejamiento evento--elemento durante 2021--2023— que de otro modo habría comprometido silenciosamente la calidad de los resultados. La delimitación consecuente del dataset de modelado a 2024--2025 constituye, en sí misma, una métrica de integridad aplicada: se prefirió un panel más corto pero confiable a uno más extenso con sesgo no controlado.

\section{Conclusión asociada al objetivo específico 2}

Las variables operativas, técnicas y climáticas más relevantes para predecir SAIDI y SAIFI, identificadas de forma convergente por la importancia de variables del Random Forest y por los coeficientes de los modelos espaciales SAR/SEM, corresponden principalmente a características de infraestructura y exposición territorial: el número de elementos de red y la longitud de red del cuadrante, la antigüedad del activo más viejo, la composición de estrato y de clientes no residenciales, y en menor medida la exposición sísmica y la vegetación (NDVI). En contraste, la señal climática contemporánea (temperatura, humedad, viento y precipitación de la misma semana) resultó comparativamente débil una vez controlada la redundancia entre variables climáticas correlacionadas (VIF reducido de un máximo de $43{.}92$ a $7{.}21$). Este hallazgo es relevante porque contradice la expectativa inicial de que el clima sería el principal predictor, y en cambio señala que la infraestructura de red —una variable sobre la que el operador tiene margen de gestión directa— es el factor más informativo disponible \emph{ex ante}.

\section{Conclusión asociada al objetivo específico 3}

Se evaluó el desempeño de seis métodos —cuatro de panel semanal (Pooled OLS, efectos fijos de cuadrante, efectos fijos de cuadrante y semana, Random Forest) y dos modelos espaciales formales (SAR, SEM)— para predecir SAIDI y SAIFI por cuadrante, considerando tanto la naturaleza de los datos como la interpretabilidad requerida por los gestores. Random Forest resultó el modelo de pronóstico semanal con menor error para SAIDI (RMSE $=0{.}1195$), mientras que el modelo de efectos fijos de cuadrante fue superior para SAIFI (RMSE $=0{.}0026$), evidenciando que no existe un método único óptimo para ambos indicadores. Para la explicación estructural, los modelos espaciales SAR (SAIDI) y SEM (SAIFI) resultaron preferidos por AIC sobre el corte transversal ponderado por cuadrante, con una capacidad explicativa alta cuantificada mediante su regresión lineal equivalente ($R^2 = 0{.}901$ y $R^2 = 0{.}837$, respectivamente), verificada mediante pruebas de robustez que descartan que dicha capacidad sea un artefacto de colinealidad. Estos resultados permiten recomendar de forma diferenciada, según el propósito operativo: Random Forest o efectos fijos de cuadrante para pronóstico de corto plazo, y SAR/SEM para el diagnóstico territorial estructural de la calidad del servicio.
